\YL{[please fill in here a short discussion about the BPZ formalism of CFT]}
\YL{[follow Mnev's lecture notes on CFT]}

The above three sections we explicitly and implicitly made some assumptions apart from the assumption of conformal symmetry :
\begin{itemize}
  \item A CFT should contain Primary Fields and their descendants only
  \item General CFT should contain a E-M Tensor with a suitable OPE
  \item Conformal Ward identity holds beyond primary fields
  \item A global conformal invariant vacuum $ |0\rangle $ 
  \item Hilbert Space is spanned by the highest weight representations of the Virasoro Algebra
  \item ...
\end{itemize}
These assumptions are really messy. However, they made the final structure of CFT beautiful and clear. Thus, we generally won't use these explicit and implicit assumptions in the above sections to define a CFT (or we will fall into mess again), but we will make the clear results (OPE Structure, Virasoro Algebra Representation Spectrum , state operator correspondence...) as a defining feature of a CFT. These feature gives a abstract but clear definition of a CFT, this work is first done by Belavin, Polyakov and Zamolodchikov in their famous BPZ paper thus I'll here give a brief introduction to BPZ formalism of CFT. 

\subsection{BPZ Axioms of CFT}




\subsection{}




